\documentclass[12pt, a4paper]{article}

% --- Codificação e língua ---
\usepackage[utf8]{inputenc}
\usepackage[T1]{fontenc}
\usepackage[brazil]{babel}

% --- Matemática ---
\usepackage{amsmath}
\usepackage{amssymb}

% --- Tabelas e figuras ---
\usepackage{booktabs}
\usepackage{graphicx}
\usepackage{float}
\usepackage{caption}
\usepackage{subcaption}

% --- Caminho das figuras ---
\graphicspath{{C:/Users/felip/OneDrive/Documentos/tcc/Felipe_Dornelles_tcc/forecasting_inflation/results/resultsk20lvar085kfold3/}}

% --- Margens ---
\usepackage[top=2.5cm, bottom=2.5cm,
            left=3cm,  right=2cm]{geometry}

% --- Espaçamento ---
\usepackage{setspace}
\onehalfspacing

% --- Referências ---
\usepackage{natbib}
\bibliographystyle{apalike}

% --- Hiperlinks ---
\usepackage[colorlinks=true, linkcolor=blue,
            citecolor=blue, urlcolor=blue]{hyperref}

% ============================================================
\begin{document}

\title{\textbf{Previsão de Inflação com Parâmetros Variantes no Tempo:\\
       Modelo 2SRR aplicado ao CPIAUCSL}\\[0.5em]
       \large Hiperparâmetros: $K_{\max}=20$,
       Var.\ expl.\ $= 85\%$, $k$-fold $= 3$}
\author{Felipe Dornelles}
\date{\today}
\maketitle
\tableofcontents
\newpage

% ============================================================
\section{Modelo 2SRR: Parâmetros Variantes no Tempo como Regressões Ridge}
\label{sec:2srr}

\subsection{Motivação}

Os modelos de previsão com parâmetros fixos assumem que a relação entre
preditores e inflação permanece estável ao longo do tempo. Contudo,
eventos como a crise financeira de 2008 e o choque inflacionário
pós-pandemia de 2021--2022 sugerem que essa relação é fundamentalmente
instável. \citet{Coulombe2021} demonstra que modelos com parâmetros
variantes no tempo (\textit{time-varying parameters}, TVP) podem ser
interpretados como casos especiais de regressão Ridge com penalidade
sobre as mudanças dos coeficientes, o que abre caminho para estimação
eficiente em grandes painéis de variáveis.

\subsection{Especificação do Modelo}

Seja $y_{t+h}$ a inflação acumulada em $h$ meses à frente e
$\mathbf{x}_t \in \mathbb{R}^N$ o vetor de $N$ preditores
macroeconômicos disponíveis em $t$. O modelo 2SRR
(\textit{Two-Step Ridge Regression}) é estimado em dois estágios:

\paragraph{Estágio 1 — Redução de dimensionalidade via PCA.}
Seja $\mathbf{F}_t \in \mathbb{R}^K$ o vetor dos $K$ primeiros
componentes principais de $\mathbf{x}_t$, selecionados de modo a
explicar ao menos 85\% da variância total ($K \leq 20$). A
transformação é:
\[
  \mathbf{F}_t = \mathbf{P}^\top \tilde{\mathbf{x}}_t
\]
onde $\mathbf{P}$ é a matriz de autovetores e $\tilde{\mathbf{x}}_t$
é a versão padronizada de $\mathbf{x}_t$.

\paragraph{Estágio 2 — Regressão Ridge com parâmetros variantes.}
O vetor de coeficientes $\boldsymbol{\beta}(t)$ é estimado via Ridge:
\[
  \hat{\boldsymbol{\beta}}(\lambda) =
  \arg\min_{\boldsymbol{\beta}}
  \sum_{s=1}^{T} \left(y_{s+h} - \mathbf{F}_s^\top \boldsymbol{\beta}_s
  \right)^2
  + \lambda \sum_{s=2}^{T}
    \|\boldsymbol{\beta}_s - \boldsymbol{\beta}_{s-1}\|^2
\]
O parâmetro de penalidade $\lambda$ é selecionado por validação
cruzada $k$-fold ($k=3$) em cada janela de estimação. Como mostrado
por \citet{Coulombe2021}, essa formulação é algebricamente equivalente
a um modelo TVP-Ridge, onde $\lambda$ controla a velocidade de
variação dos coeficientes no tempo.

\subsection{Implementação}

O modelo foi implementado na função \texttt{run2srr()} em R, com os
seguintes hiperparâmetros:

\begin{itemize}
  \item Número máximo de componentes principais: $K_{\max} = 20$
  \item Variância explicada mínima: $85\%$
  \item Validação cruzada: $k\text{-fold} = 3$
  \item Horizontes estimados: $h \in \{1, 2, \ldots, 12\}$ meses
  \item Janelas fora da amostra: 180 (julho de 2010 a junho de 2025)
\end{itemize}

Os coeficientes $\hat{\boldsymbol{\beta}}(t)$ foram extraídos para
cada janela e horizonte, totalizando $180 \times 21 = 3.780$
observações por horizonte.

% ============================================================
\section{Resultados}
\label{sec:resultados}

\subsection{Acurácia Preditiva: RMSFE relativo ao Random Walk}

A métrica principal de avaliação é o RMSFE
(\textit{Root Mean Squared Forecast Error}) relativo ao benchmark
de Random Walk (RW):
\[
  \text{RMSFE}_{\text{rel}}(h) =
  \frac%
    {\sqrt{\dfrac{1}{T}\displaystyle\sum_t
      \bigl(f_{t,h} - y_{t+h}\bigr)^2}}%
    {\sqrt{\dfrac{1}{T}\displaystyle\sum_t
      \bigl(y_t    - y_{t+h}\bigr)^2}}
\]
Valores abaixo de 1 indicam que o modelo supera o Random Walk.

\begin{table}[H]
\centering
\caption{RMSFE relativo ao Random Walk --- principais modelos.
  Negrito indica o melhor resultado por linha.
  Valores $< 1$ indicam ganho sobre o Random Walk.}
\label{tab:rmsfe}
\begin{tabular}{lrrrrr}
\toprule
Horizonte & \textbf{2SRR} & Ridge  & LASSO  & RF     & Factor \\
\midrule
$h=1$  & 0{,}9765 & 1{,}1636 & 1{,}1702 & 1{,}1528 & 1{,}2803 \\
$h=2$  & \textbf{0{,}8720} & 0{,}9818 & 0{,}8808 & 0{,}8895 & 0{,}9773 \\
$h=3$  & 0{,}9445 & 0{,}9222 & \textbf{0{,}8278} & 0{,}8690 & 0{,}9575 \\
$h=4$  & 0{,}9174 & 0{,}9598 & \textbf{0{,}8777} & 0{,}8945 & 0{,}9252 \\
$h=5$  & \textbf{0{,}8429} & 0{,}9490 & 0{,}8650 & 0{,}8540 & 0{,}9641 \\
$h=6$  & \textbf{0{,}8042} & 0{,}9656 & 0{,}8690 & 0{,}9065 & 0{,}9547 \\
$h=7$  & 0{,}8384 & 0{,}9136 & \textbf{0{,}8363} & 0{,}8683 & 0{,}9543 \\
$h=8$  & \textbf{0{,}8570} & 1{,}0514 & 0{,}9216 & 0{,}8856 & 1{,}0794 \\
$h=9$  & 0{,}8314 & 0{,}9011 & 0{,}7842 & \textbf{0{,}7693} & 0{,}8998 \\
$h=10$ & \textbf{0{,}8782} & 1{,}0153 & 0{,}9127 & 0{,}9141 & 0{,}9868 \\
$h=11$ & \textbf{0{,}8652} & 0{,}9787 & 0{,}8675 & 0{,}8900 & 0{,}9451 \\
$h=12$ & \textbf{0{,}8521} & 1{,}0435 & 0{,}9437 & 0{,}9782 & 1{,}0278 \\
\midrule
acc3   & \textbf{0{,}6689} & 0{,}7344 & 0{,}6865 & 0{,}6643 & 0{,}8427 \\
acc6   & 0{,}5349 & \textbf{0{,}5114} & 0{,}4775 & 0{,}4793 & 0{,}6370 \\
acc12  & 0{,}3327 & 0{,}2556 & \textbf{0{,}2335} & 0{,}3552 & 0{,}7787 \\
\bottomrule
\end{tabular}
\end{table}

O 2SRR supera o Random Walk em 11 dos 12 horizontes individuais,
com ganhos expressivos nos horizontes de médio prazo ($h = 5$ a
$h = 9$). Nos acumulados, o modelo apresenta RMSFE relativo de
$0{,}669$ no acc3 --- o menor entre todos os modelos avaliados ---,
indicando forte capacidade preditiva para inflação trimestral
acumulada. O desempenho inferior no acc12 frente ao Ridge e ao LASSO
é esperado: o \textit{shrinkage} fixo desses modelos favorece
horizontes muito longos, ao passo que o 2SRR adapta os coeficientes
ao longo do tempo, vantagem concentrada nos horizontes intermediários.

\begin{figure}[H]
  \centering
  \includegraphics[width=\textwidth]{fig1_rmsfe_todos.png}
  \caption{RMSFE relativo ao Random Walk para todos os modelos,
    por horizonte ($h = 1$ a $h = 12$). A linha tracejada representa
    $\text{RW} = 1$. O 2SRR (linha vermelha em destaque) mantém-se
    abaixo de 1 em todos os horizontes a partir de $h = 2$, sendo
    o modelo de trajetória mais estável do conjunto avaliado.}
  \label{fig:rmsfe_todos}
\end{figure}

\begin{figure}[H]
  \centering
  \includegraphics[width=\textwidth]{fig2_2srr_vs_outros.png}
  \caption{Comparação do RMSFE relativo ao RW entre 2SRR, Ridge,
    LASSO e Random Forest, para $h = 1$ a $h = 12$ e acumulados.
    O 2SRR apresenta a trajetória mais estável, sem nenhum horizonte
    acima de 1 a partir de $h = 2$, ao contrário do Ridge, que supera
    o RW em $h = 8$, $h = 10$ e $h = 12$.}
  \label{fig:rmsfe_comparativo}
\end{figure}

\subsection{Parâmetros Variantes no Tempo}

Uma contribuição distintiva do 2SRR é a possibilidade de recuperar
e interpretar a trajetória temporal dos coeficientes
$\hat{\boldsymbol{\beta}}(t)$. A Figura~\ref{fig:betas_tv} apresenta
a evolução dos coeficientes das 10 variáveis de maior variância,
para $h = 1$.

\begin{figure}[H]
  \centering
  \includegraphics[width=\textwidth]{fig_betas_2SRR_tempo.png}
  \caption{Evolução de $\hat{\beta}(t)$ do 2SRR para $h = 1$,
    top 10 variáveis por variância dos coeficientes (2010--2025).
    Nota-se aumento expressivo da volatilidade no período 2021--2025,
    coincidindo com o choque inflacionário pós-pandemia.}
  \label{fig:betas_tv}
\end{figure}

Três padrões se destacam:

\begin{enumerate}
  \item \textbf{Estabilidade relativa (2010--2019):} coeficientes
    oscilam em torno de zero com amplitude moderada, consistente
    com inflação baixa e estável nos EUA.
  \item \textbf{Ruptura estrutural (2020--2021):} queda abrupta
    nos betas de variáveis industriais (\texttt{IPDMAT},
    \texttt{IPFINAL}), refletindo o colapso da atividade real.
  \item \textbf{Reajuste inflacionário (2021--2025):} variáveis de
    manufatura (\texttt{IPDCONGD}, \texttt{IPCONGD},
    \texttt{IPMANSICS}) passam a apresentar coeficientes de grande
    magnitude, capturando o papel da oferta industrial no episódio
    inflacionário recente.
\end{enumerate}

Esse comportamento confirma a hipótese central de
\citet{Coulombe2021}: a relação entre preditores e inflação é
estruturalmente instável, e modelos com parâmetros fixos
subestimam sistematicamente essa heterogeneidade temporal.

% ============================================================
\bibliography{references}

\end{document}